% !TeX spellcheck = pt_BR
\documentclass[12pt,portuguese]{article}

%----------------------------------------------------------------------
% PACOTES ESSENCIAIS E TRADUÇÃO
%----------------------------------------------------------------------
\usepackage[brazilian]{babel}     % Suporte à língua portuguesa
\usepackage[utf8]{inputenc}       % Codificação de entrada
\usepackage[T1]{fontenc}          % Codificação da fonte
\usepackage{lmodern}              % Fonte escalável

%----------------------------------------------------------------------
% PACOTES MATEMÁTICOS
%----------------------------------------------------------------------
\usepackage{amsmath, amsthm, amssymb, amstext, mathtools}
\usepackage{mathrsfs}             % Para fontes como \mathscr
\usepackage{dsfont}               % Para o indicador de conjunto \mathds{1}
\usepackage[all]{xy}              % Para diagramas comutativos
\usepackage{imakeidx}             % Criação de Lista de Símbolos / Índice

% Desigualdades sempre com símbolos inclinados
\let\leq\leqslant
\let\le\leqslant
\let\ge\geqslant
\let\geq\geqslant

%----------------------------------------------------------------------
% FORMATAÇÃO, LAYOUT E CORES
%----------------------------------------------------------------------
\usepackage{geometry}
\usepackage{varwidth}
\usepackage{graphicx}
\usepackage[pdftex]{xcolor}
\usepackage{titlesec}
\usepackage{appendix}
\usepackage{enumitem}
\usepackage{float}
\usepackage{booktabs}
\usepackage{array}
\usepackage{colortbl}
\usepackage{fontawesome5}
\usepackage{lucide-icons}
\usepackage{caption}
\usepackage{stackengine}
\captionsetup{width=0.8\textwidth}

% Definição de Cores
\definecolor{teal}{RGB}{15,118,110}
\definecolor{text}{RGB}{28,25,23}
\definecolor{muted}{RGB}{87,83,78}
\definecolor{card}{RGB}{255,255,255}
\definecolor{border}{RGB}{231,229,228}

% Definição de margens e espaçamento
\setlength{\topmargin}{-0.7in}
\setlength{\oddsidemargin}{0in}
\setlength{\textheight}{9.55in}
\setlength{\textwidth}{6.5in}
\setlength{\baselineskip}{12mm}

% Formatação dos títulos de seção
\titleformat{\section}{\Large\bfseries}{\thesection.}{0.5em}{}
\titleformat{\subsection}{\large\bfseries}{\thesubsection.}{0.5em}{}
\titleformat{\subsubsection}{\normalfont\itshape}{\thesubsubsection.}{0.6em}{}

%----------------------------------------------------------------------
% TIKZ E GRÁFICOS
%----------------------------------------------------------------------
\usepackage{tikz}
\usepackage{pgfplots}
\pgfplotsset{compat=1.17}
\usetikzlibrary{positioning,calc,fit,nfold,decorations.pathreplacing,calligraphy,tikzmark,patterns,arrows.meta,bending,hobby}

%----------------------------------------------------------------------
% REFERÊNCIAS, LINKS E METADADOS DO PDF
%----------------------------------------------------------------------
\usepackage[pdftex, colorlinks, linkcolor={red!50!black}, 
citecolor={red}, urlcolor={blue!80!black}]{hyperref}

\hypersetup{
pdftitle={Funções de Distribuição},
pdfauthor={L. Cioletti},
pdfsubject={Funções de distribuição, variáveis aleatórias contínuas, funções densidade de probabilidade, densidades usuais},
pdfkeywords={função de distribuição, variável aleatória contínua, função densidade de probabilidade, densidade, distribuição uniforme, exponencial, normal, gama, beta, MAT-UnB}
}

% --- PACOTE E COMANDO PARA A LICENÇA CREATIVE COMMONS ---
\usepackage[type={CC},modifier={by-nc-sa},version={4.0}]{doclicense}
\newcommand\blfootnote[1]{%
\begingroup
\renewcommand\thefootnote{}\footnote{#1}%
\addtocounter{footnote}{-1}%
\endgroup
}

%----------------------------------------------------------------------
% ALIASCNT E AMBIENTES DE TEOREMA (Cleveref)
%----------------------------------------------------------------------
\usepackage{aliascnt}
\theoremstyle{definition}
\newtheorem{theorem}{Teorema}

\newaliascnt{definition}{theorem}
\newtheorem{definition}[definition]{Definição}
\aliascntresetthe{definition}

\newaliascnt{example}{theorem}
\newtheorem{example}[example]{Exemplo}
\aliascntresetthe{example}

\newaliascnt{proposition}{theorem}
\newtheorem{proposition}[proposition]{Proposição}
\aliascntresetthe{proposition}

\newaliascnt{corollary}{theorem}
\newtheorem{corollary}[corollary]{Corolário}
\aliascntresetthe{corollary}

\newaliascnt{lemma}{theorem}
\newtheorem{lemma}[lemma]{Lema}
\aliascntresetthe{lemma}

\newaliascnt{remark}{theorem}
\newtheorem{remark}[remark]{Observação}
\aliascntresetthe{remark}

\newtheoremstyle{plainexercise}
  {3pt}{3pt}{\normalfont}{}{\bfseries}{.}{.5em}
  {\thmnumber{#2}\thmnote{ #3}}
\theoremstyle{plainexercise}
\newtheorem{exercise}{Exercício}

% Versão marcada dos exercícios: a estrela fica pendurada na margem esquerda,
% como nas notas anteriores.
\newtheoremstyle{starredexercise}
  {3pt}{3pt}{\normalfont}{}{\bfseries}{.}{.5em}
  {\llap{\textasteriskcentered\,}\thmnumber{#2}\thmnote{ #3}}
\theoremstyle{starredexercise}
\newtheorem{starredexercise}[exercise]{Exercício}
\theoremstyle{definition}

% Ambientes não numerados usados no texto
\newtheorem*{property}{Propriedade}
\newtheorem*{solution}{Solução}

% Se um enumerate for o primeiro conteúdo de um teorema/exercício (cabeçalho
% ainda pendente, \@inlabel true), o \leavevmode dispara o \everypar do amsthm
% que coloca o cabeçalho na própria linha, antes da lista iniciar.
\makeatletter
\setlist[enumerate]{before={\if@inlabel\leavevmode\fi}}
\makeatother

% Nomes em português para os ambientes
\addto\captionsbrazil{\renewcommand{\proofname}{\normalfont\textbf{Prova}}}
\renewcommand{\qedsymbol}{$\blacksquare$}

% Cleveref (deve vir depois de hyperref e amsthm)
\usepackage[brazilian,nameinlink]{cleveref}

\crefname{theorem}{teorema}{teoremas}
\Crefname{theorem}{Teorema}{Teoremas}
\crefname{definition}{definição}{definições}
\Crefname{definition}{Definição}{Definições}
\crefname{lemma}{lema}{lemas}
\Crefname{lemma}{Lema}{Lemas}
\crefname{proposition}{proposição}{proposições}
\Crefname{proposition}{Proposição}{Proposições}
\crefname{corollary}{corolário}{corolários}
\Crefname{corollary}{Corolário}{Corolários}
\crefname{example}{exemplo}{exemplos}
\Crefname{example}{Exemplo}{Exemplos}
\crefname{exercise}{exercício}{exercícios}
\Crefname{exercise}{Exercício}{Exercícios}
\crefname{remark}{observação}{observações}
\Crefname{remark}{Observação}{Observações}
\crefname{section}{seção}{seções}
\Crefname{section}{Seção}{Seções}

%----------------------------------------------------------------------
% MACROS ADICIONAIS
%----------------------------------------------------------------------
\let\oldemptyset\emptyset
\let\emptyset\varnothing
\graphicspath{{./}{../arquivos-aux/}}

% Seno na convenção em português
\renewcommand{\sin}{\operatorname{sen}}

%----------------------------------------------------------------------
% DOCUMENTO PRINCIPAL
%----------------------------------------------------------------------
\begin{document}
% --- CABEÇALHO COM O LOGO DA UnB ---
\begin{figure}[H]
\includegraphics{logo-mat-unb.pdf}
\end{figure}

\vspace{0.1cm}

% --- TÍTULO, AUTOR E DATA ---
\begin{center}
{\LARGE\bf Funções de Distribuição}
\vspace{1cm}

{\large \bf L. Cioletti}
\vspace{0.3cm}

{\footnotesize Setembro de 2026}
\end{center}

% --- RESUMO (ABSTRACT) ---
\begin{abstract}
Estas notas dão continuidade ao estudo das variáveis aleatórias iniciado nas
notas sobre variáveis aleatórias discretas. Introduzimos a noção geral de
variável aleatória e a função de distribuição, estudamos suas propriedades
gerais e apresentamos os primeiros exemplos. Em seguida definimos as variáveis
aleatórias contínuas por meio de suas funções densidade de probabilidade e
mostramos como a densidade se relaciona com a função de distribuição.
Reunimos também as densidades usuais, a saber, a uniforme, a exponencial, a
normal, a de Cauchy, a gama, a beta e a qui-quadrado, verificando as suas
normalizações e a relação entre as funções gama e beta. O texto segue de perto
as Seções 5.1--5.4 do Capítulo 5 de Grimmett e Welsh \cite{GrimmettWelsh2014}
e conta com uma lista de exercícios ao final.
\end{abstract}

% --- LICENÇA NO RODAPÉ ---
\blfootnote{%
\doclicenseImage[imagewidth=2.6em]\ 
\begingroup
\hypersetup{urlcolor=darkgray}
\href{\doclicenseURL}{%
\doclicenseName - Attribution-NonCommercial-ShareAlike 4.0 International%
}%
\endgroup
}

%----------------------------------------------------------------------
% CONTEÚDO
%----------------------------------------------------------------------

\section{Funções de Distribuição}
\label{sec:funcoes-distribuicao}

Até aqui trabalhamos com variáveis aleatórias discretas, que assumem apenas
uma quantidade enumerável de valores. Essa condição é restritiva demais para
muitas situações e, por isso, vamos adotar uma definição mais ampla.

\begin{definition}[Variável aleatória]
\label{def:variavel-aleatoria}
Seja $(\Omega,\mathcal{F},\mathbb{P})$ um espaço de probabilidade. Dizemos
que $X$ é uma \emph{variável aleatória} se $X$ é uma função
$X:\Omega\to\mathbb{R}$ tal que
\begin{equation}
\label{eq:variavel-aleatoria}
  \{\omega\in\Omega : X(\omega)\leqslant x\}\in\mathcal{F}
  \quad\text{para todo } x\in\mathbb{R}.
\end{equation}
\end{definition}

Quando \eqref{eq:variavel-aleatoria} vale, dizemos que $X$ é
$\mathcal{F}$-mensurável. Daqui em diante, abreviamos o evento
$\{\omega\in\Omega : X(\omega)\leqslant x\}$ pela expressão mais simples
$\{X\leqslant x\}$.

Exigimos que as variáveis aleatórias satisfaçam \eqref{eq:variavel-aleatoria}
pela mesma razão pela qual exigimos uma condição análoga na definição de
variável aleatória discreta. Estamos interessados nos valores assumidos por
$X$ e nas probabilidades associadas a esses valores. Uma maneira conveniente
de estudá-los é fixar $x\in\mathbb{R}$ e calcular a probabilidade de que $X$
assuma um valor no intervalo $(-\infty,x]$. Essa probabilidade só está
definida quando a pré-imagem
$X^{-1}\big((-\infty,x]\big)=\{X\leqslant x\}$ pertence ao espaço de eventos
$\mathcal{F}$, e é por isso que postulamos que essa condição vale para todo
$x\in\mathbb{R}$.

Observe que toda variável aleatória discreta é uma variável aleatória. De
fato, se $X$ é discreta, então
\[
  \{X\leqslant x\}
  = \bigcup_{y\in\operatorname{Im}X :\, y\leqslant x} \{X=y\},
\]
que é uma união enumerável de eventos em $\mathcal{F}$ e, portanto, também
pertence a $\mathcal{F}$.

Enquanto as variáveis aleatórias discretas eram estudadas por meio de suas
funções de probabilidade, as variáveis aleatórias no sentido mais amplo são
estudadas por meio de suas funções de distribuição, definidas a seguir.

\begin{definition}[Função de distribuição]
\label{def:funcao-distribuicao}
Seja $X$ uma variável aleatória em $(\Omega,\mathcal{F},\mathbb{P})$. A
\emph{função de distribuição} de $X$ é a função
$F_X:\mathbb{R}\to[0,1]$ definida por
\begin{equation}
\label{eq:def-funcao-distribuicao}
  F_X(x)=\mathbb{P}(X\leqslant x).
\end{equation}
\end{definition}

A função de distribuição de $X$ também é chamada, em alguns textos, de função
de distribuição acumulada de $X$.

\begin{example}[Função de distribuição de uma variável discreta]
\label{ex:fd-discreta}
Suponha que $X$ seja uma variável aleatória discreta que assume valores
inteiros não negativos, com função de probabilidade
\[
  \mathbb{P}(X=k)=p_k \quad\text{para } k=0,1,2,\dots.
\]
Para $x\in\mathbb{R}$, temos que $X\leqslant x$ se, e somente se, $X$ assume
um dos valores $0,1,2,\dots,\lfloor x\rfloor$, onde $\lfloor x\rfloor$ denota
o maior inteiro que não supera $x$. Logo
\[
  F_X(x)=
  \begin{cases}
    0, & \text{se } x<0,\\[0.2cm]
    p_0+p_1+\dots+p_{\lfloor x\rfloor}, & \text{se } x\geqslant 0.
  \end{cases}
\]
Um esboço dessa função, no caso em que $X$ assume valores no conjunto
$\{0,1,2,3,4,5\}$, é apresentado na \Cref{fig:fd-discreta}.
\end{example}

\begin{figure}[H]
\centering
\begin{tikzpicture}[xscale=1.0,yscale=2.4]
  % eixos
  \draw[->] (-2.4,0) -- (6.6,0) node[below right] {$x$};
  \draw[->] (0,-0.06) -- (0,1.18) node[above left] {$F_X(x)$};
  % patamares, contínua à direita
  \draw[teal,thick] (-2.4,0) -- (0,0);
  \draw[teal,thick] (0,0) -- (0,0.10);
  \draw[teal,thick] (0,0.10) -- (1,0.10);
  \draw[teal,thick] (1,0.10) -- (1,0.28);
  \draw[teal,thick] (1,0.28) -- (2,0.28);
  \draw[teal,thick] (2,0.28) -- (2,0.46);
  \draw[teal,thick] (2,0.46) -- (3,0.46);
  \draw[teal,thick] (3,0.46) -- (3,0.64);
  \draw[teal,thick] (3,0.64) -- (4,0.64);
  \draw[teal,thick] (4,0.64) -- (4,0.80);
  \draw[teal,thick] (4,0.80) -- (5,0.80);
  \draw[teal,thick] (5,0.80) -- (5,1.00);
  \draw[teal,thick] (5,1.00) -- (6.4,1.00);
  % pontos incluídos
  \foreach \p in {(0,0.10),(1,0.28),(2,0.46),(3,0.64),(4,0.80),(5,1.00)}
    \fill[teal] \p circle (1.6pt);
  % pontos excluídos
  \foreach \p in {(0,0),(1,0.10),(2,0.28),(3,0.46),(4,0.64),(5,0.80)}
    \draw[teal,fill=white,thick] \p circle (1.6pt);
  % marcas no eixo horizontal
  \foreach \x in {0,1,2,3,4,5} {
    \draw (\x,0.02) -- (\x,-0.02);
    \node[below] at (\x,-0.03) {\small $\x$};
  }
  % marca no eixo vertical
  \draw (-0.05,1.00) -- (0.05,1.00);
  \node[left] at (-0.07,1.00) {\small $1$};
\end{tikzpicture}
\caption{A função de distribuição de uma variável aleatória que assume
valores no conjunto $\{0,1,2,3,4,5\}$. Os degraus são crescentes e a função é
contínua à direita, mas descontínua à esquerda em cada ponto do suporte.}
\label{fig:fd-discreta}
\end{figure}

A função de distribuição $F_X$ possui várias propriedades gerais e
elementares, reunidas a seguir.

\begin{proposition}[Propriedades da função de distribuição]
\label{prop:propriedades-funcao-distribuicao}
Seja $X$ uma variável aleatória em $(\Omega,\mathcal{F},\mathbb{P})$ e seja
$F_X$ a sua função de distribuição. Então:

\begin{enumerate}[label=\textup{(\roman*)}]
  \item $F_X$ é monótona não decrescente, isto é,
    \begin{equation}
    \label{eq:fd-monotona}
      F_X(x)\leqslant F_X(y)\quad\text{se } x\leqslant y;
    \end{equation}
  \item os limites nas extremidades são
    \begin{equation}
    \label{eq:fd-limites}
      \lim_{x\to-\infty}F_X(x)=0
      \qquad\text{e}\qquad
      \lim_{x\to+\infty}F_X(x)=1;
    \end{equation}
  \item $F_X$ é contínua à direita, isto é,
    \begin{equation}
    \label{eq:fd-continua-direita}
      F_X(x+\epsilon)\to F_X(x)\quad\text{quando } \epsilon\downarrow 0,
    \end{equation}
    onde o limite é tomado com $\epsilon$ tendendo a $0$ apenas por valores
    positivos.
\end{enumerate}
\end{proposition}

A monotonicidade expressa em \eqref{eq:fd-monotona} decorre da inclusão
\[
  \{X\leqslant x\}\subseteq\{X\leqslant y\}
\]
sempre que $x\leqslant y$, pois se $X$ assume um valor que não excede $x$,
então esse valor também não excede $y$.

As demais propriedades elementares de $F_X$ dizem respeito ao seu
comportamento quando $x$ se aproxima de $-\infty$ ou de $+\infty$.
Intuitivamente é claro que
\[
  F_X(x)\to 0 \quad\text{quando } x\to-\infty,
  \qquad\text{e}\qquad
  F_X(x)\to 1 \quad\text{quando } x\to+\infty,
\]
uma vez que, no primeiro caso, à medida que $x$ diminui, o evento de que $X$
seja menor do que $x$ torna-se cada vez menos provável; no segundo caso, à
medida que $x$ cresce, esse evento torna-se cada vez mais provável. No plano
intuitivo, as igualdades em \eqref{eq:fd-limites} são evidentes, pois
lembram as afirmações
\[
  \mathbb{P}(X\leqslant-\infty)=0,
  \qquad
  \mathbb{P}(X\leqslant+\infty)=1,
\]
mas a verificação formal de \eqref{eq:fd-limites} depende da continuidade da
medida de probabilidade $\mathbb{P}$.

Para tornar essa verificação formal, considere os eventos
$A_n=\{X\leqslant-n\}$. Eles formam uma sequência decrescente e
$\bigcap_{n\geqslant1}A_n=\varnothing$, pois nenhum número real é menor ou
igual a $-n$ para todo $n$. Pela continuidade por cima da probabilidade,
\[
  F_X(-n)=\mathbb{P}(A_n)\longrightarrow 0.
\]
Analogamente, se $B_n=\{X\leqslant n\}$, então $B_n$ cresce para $\Omega$ e,
pela continuidade por baixo,
\[
  F_X(n)=\mathbb{P}(B_n)\longrightarrow 1.
\]
Como $F_X$ é monótona, esses dois limites ao longo dos inteiros implicam os
limites em \eqref{eq:fd-limites}.

A continuidade da medida de probabilidade também é necessária para
demonstrar a continuidade à direita em \eqref{eq:fd-continua-direita}. De
fato, se $(x_n)_{n\geqslant1}$ é uma sequência decrescente com $x_n\to x$,
então os eventos $\{X\leqslant x_n\}$ decrescem para $\{X\leqslant x\}$ e a
continuidade por cima de $\mathbb{P}$ dá
\[
  F_X(x_n)=\mathbb{P}(X\leqslant x_n)\to\mathbb{P}(X\leqslant x)=F_X(x).
\]
Observando a \Cref{fig:fd-discreta}, vemos que as funções de distribuição não
precisam ser contínuas à esquerda: em cada ponto do suporte ocorre um salto.

\begin{remark}
\label{rem:caracterizacao-fd}
As três propriedades reunidas na
\Cref{prop:propriedades-funcao-distribuicao}, isto é, a monotonicidade
\eqref{eq:fd-monotona}, os limites \eqref{eq:fd-limites} e a continuidade à
direita \eqref{eq:fd-continua-direita}, caracterizam completamente as funções
de distribuição. Isso significa que, se $F$ é uma função que satisfaz essas
três condições, então existem um espaço de probabilidade e uma variável
aleatória $X$ nesse espaço cuja função de distribuição é $F$. A prova desse
fato é omitida nestas notas, mas convém registrá-lo, pois em muitas situações
ele nos permite evitar o trabalho de exibir explicitamente o espaço de
probabilidade e a variável aleatória.
\end{remark}

Antes de apresentarmos exemplos de funções de distribuição, vejamos uma
última propriedade. A probabilidade $F_X(x)=\mathbb{P}(X\leqslant x)$ é a
probabilidade de que $X$ assuma um valor no intervalo ilimitado
$(-\infty,x]$. Para encontrarmos a probabilidade de que $X$ assuma um valor
em um intervalo limitado $(a,b]$, procedemos da seguinte maneira. Para
$a<b$,
\begin{equation}
\label{eq:intervalo-fd}
  \mathbb{P}(a<X\leqslant b)=F_X(b)-F_X(a).
\end{equation}

De fato,
\[
  \begin{aligned}
    \mathbb{P}(a<X\leqslant b)
    &= \mathbb{P}\big(\{X\leqslant b\}\setminus\{X\leqslant a\}\big)\\
    &= \mathbb{P}(X\leqslant b)-\mathbb{P}(X\leqslant a),
  \end{aligned}
\]
pois o evento $\{X\leqslant a\}$ está contido no evento $\{X\leqslant b\}$.

\section{Exemplos de Funções de Distribuição}\label{sec:exemplos-fd}

Na seção anterior introduzimos a definição de função de distribuição e reunimos
as suas propriedades gerais na \Cref{prop:propriedades-funcao-distribuicao}.
O \Cref{ex:fd-discreta} foi o nosso primeiro exemplo: a função de distribuição
de uma variável aleatória discreta, constante por partes e com saltos
exatamente nos pontos com probabilidade positiva. Vale a pena observar as características gerais
daquela função: ela é não decrescente, contínua à direita, tende a $0$ quando
$x \to -\infty$ e tende a $1$ quando $x \to \infty$. Existem, contudo, funções
de distribuição que contrastam de modo acentuado com aquele exemplo, por serem
contínuas em toda a reta. O próximo exemplo é de uma dessas funções.

\begin{example}[Distribuição uniforme]\label{ex:uniforme-fd}
  Sejam $a, b \in \mathbb{R}$ com $a < b$. Considere a função
  \[
    F(x) = \begin{cases}
      0, & \text{se } x < a, \\[0.2cm]
      \dfrac{x-a}{b-a}, & \text{se } a \leqslant x \leqslant b, \\[0.2cm]
      1, & \text{se } x > b,
    \end{cases}
  \]
  cujo gráfico está esboçado na \Cref{fig:fd-uniforme}. Ela possui as
  propriedades gerais descritas na \Cref{prop:propriedades-funcao-distribuicao}
  e, portanto, é uma função de distribuição. Dizemos que uma variável aleatória
  com essa função de distribuição tem \emph{distribuição uniforme} no intervalo
  $(a,b)$; alguns autores preferem dizer que a distribuição é uniforme no
  intervalo $[a,b]$.

  O fato de chamarmos a distribuição de uniforme em $(a,b)$ ou em $[a,b]$
  não altera a lei, pois os extremos têm probabilidade zero. Os valores de
  $F$ nesses pontos, contudo, são determinados pela própria função de
  distribuição. Em particular, a fórmula define uma função de distribuição
  para quaisquer escolhas de $a < b$, e a equação \eqref{eq:intervalo-fd} se
  aplica diretamente.
\end{example}

\begin{example}[Distribuição exponencial]\label{ex:exponencial-fd}
  Seja $\lambda > 0$ e defina $F$ por
  \begin{equation}\label{eq:fd-exponencial}
    F(x) = \begin{cases}
      0, & \text{se } x \leqslant 0, \\[0.2cm]
      1 - e^{-\lambda x}, & \text{se } x > 0,
    \end{cases}
  \end{equation}
  cujo gráfico esboçamos na \Cref{fig:fd-exponencial}. Vemos que $F$ é
  claramente uma função de distribuição. Dizemos que uma variável aleatória com
  essa distribuição tem \emph{distribuição exponencial} com parâmetro
  $\lambda$.
\end{example}

As duas funções de distribuição acima são muito importantes na teoria da
probabilidade. As duas são, além disso, casos particulares de uma família
maior. Pela caracterização da \Cref{rem:caracterizacao-fd}, qualquer função
$F$ não negativa, contínua e não decrescente que satisfaz
\[
  \lim_{x \to -\infty} F(x) = 0, \qquad \lim_{x \to \infty} F(x) = 1
\]
é uma função de distribuição.

\begin{figure}[H]
\centering
\begin{tikzpicture}[xscale=1.15,yscale=1.8]
  \draw[->] (-0.6,0) -- (6.3,0) node[below right] {$x$};
  \draw[->] (0,-0.25) -- (0,1.5) node[left] {$F(x)$};
  \draw[gray,dashed] (0,1) -- (6.1,1);
  \node[left] at (0,1) {$1$};
  \node[left] at (0,0) {$0$};
  \draw (1.5,0) -- (1.5,-0.07) node[below] {$a$};
  \draw (4.5,0) -- (4.5,-0.07) node[below] {$b$};
  \draw[teal,thick] (-0.4,0) -- (1.5,0);
  \draw[teal,thick] (1.5,0) -- (4.5,1);
  \draw[teal,thick] (4.5,1) -- (6.1,1);
  \fill[teal] (1.5,0) circle (1.2pt);
  \fill[teal] (4.5,1) circle (1.2pt);
\end{tikzpicture}
\caption{Função de distribuição uniforme no intervalo $(a,b)$. A função vale
$0$ à esquerda de $a$, cresce linearmente de $0$ a $1$ em $[a,b]$ e vale $1$ à
direita de $b$.}
\label{fig:fd-uniforme}
\end{figure}

\begin{figure}[H]
\centering
\begin{tikzpicture}[xscale=0.95,yscale=1.8]
  \draw[->] (-0.6,0) -- (6.4,0) node[below right] {$x$};
  \draw[->] (0,-0.25) -- (0,1.5) node[left] {$F(x)$};
  \draw[gray,dashed] (0,1) -- (6.2,1);
  \node[left] at (0,1) {$1$};
  \node[left] at (0,0) {$0$};
  \draw[teal,thick,smooth,samples=120,domain=0:6.2]
    plot (\x,{1-exp(-0.9*\x)});
  \draw[teal,thick] (-0.4,0) -- (0,0);
  \fill[teal] (0,0) circle (1.2pt);
\end{tikzpicture}
\caption{Função de distribuição exponencial com parâmetro $\lambda$. A função é
identicamente nula em $(-\infty,0]$ e tende assintoticamente a $1$ quando
$x \to \infty$.}
\label{fig:fd-exponencial}
\end{figure}

\section{Variáveis Aleatórias Contínuas}\label{sec:va-continuas}

As variáveis aleatórias aparecem em muitas formas, mas há duas classes que
desempenham um papel central no curso:

\begin{enumerate}[label=\Roman*.]
  \item as variáveis aleatórias discretas;
  \item as variáveis aleatórias contínuas.
\end{enumerate}

Uma variável aleatória discreta assume apenas uma quantidade enumerável de
valores e sua função de distribuição tem, em geral, o aspecto de uma função
em degraus, como ilustra a \Cref{fig:fd-discreta}. No extremo oposto
encontram-se variáveis aleatórias cuja função de distribuição é bastante
suave, como nos casos uniforme e exponencial discutidos na
\Cref{sec:exemplos-fd} (veja as \Cref{fig:fd-uniforme,fig:fd-exponencial}).
É essa segunda classe que passamos a estudar e que chamaremos de
\emph{contínua}.

\begin{definition}[Variável aleatória contínua]\label{def:va-continua}
  Dizemos que uma variável aleatória $X$ é \emph{contínua} se existir uma
  função não negativa e mensurável $f_X$ tal que sua função de distribuição $F_X$ possa
  ser escrita na forma
  \begin{equation}\label{eq:densidade-integral}
    F_X(x) = \mathbb{P}(X \leqslant x) = \int_{-\infty}^{x} f_X(u)\,du
    \quad \text{para todo } x \in \mathbb{R}.
  \end{equation}
  Nesse caso, dizemos que $X$ tem \emph{função densidade de probabilidade}
  (ou simplesmente \emph{densidade}) $f_X$, abreviada por \emph{fdp}.
\end{definition}

A identidade \eqref{eq:densidade-integral} diz que a probabilidade de o evento
$\{X \leqslant x\}$ ocorrer é a área sob o gráfico de $f_X$ à esquerda de
$x$. Assim, $F_X$ acumula a densidade desde $-\infty$ até $x$, e é por isso
que $F_X$ também é chamada de função de distribuição acumulada.

\begin{remark}
  Textos mais avançados chamam as variáveis aleatórias descritas na
  \Cref{def:va-continua} de \emph{absolutamente contínuas}, reservando a
  palavra ``contínua'' para uma classe mais ampla. Neste curso não faremos
  essa distinção e usaremos ``contínua'' no sentido da definição acima.
\end{remark}

\begin{example}\label{ex:densidade-exponencial}
  Se $X$ tem distribuição exponencial de parâmetro $\lambda > 0$, então,
  como vimos no \Cref{ex:exponencial-fd},
  \[
    F_X(x) =
    \begin{cases}
      0, & \text{se } x \leqslant 0, \\[0.2cm]
      1 - e^{-\lambda x}, & \text{se } x > 0.
    \end{cases}
  \]
  Uma densidade para $X$ é obtida derivando $F_X$ nos pontos em que a
  derivada existe. Para $x > 0$,
  \[
    \frac{d}{dx}\bigl(1 - e^{-\lambda x}\bigr) = \lambda e^{-\lambda x},
  \]
  e para $x < 0$ a derivada é $0$. Logo
  \[
    f_X(x) =
    \begin{cases}
      0, & \text{se } x \leqslant 0, \\[0.2cm]
      \lambda e^{-\lambda x}, & \text{se } x > 0.
    \end{cases}
  \]
\end{example}

De modo geral, se $X$ é contínua e $F_X$ tem bom comportamento, podemos
tomar como densidade de $X$ a função
\begin{equation}\label{eq:densidade-derivada}
  f_X(x) =
  \begin{cases}
    \dfrac{d}{dx} F_X(x), & \text{se essa derivada existe em } x, \\[0.3cm]
    0, & \text{caso contrário}.
  \end{cases}
\end{equation}
Seguiremos essa receita na maioria dos exemplos. Vale registrar que há
algumas dificuldades de rigor matemático nessa passagem: nem toda função de
distribuição é derivável, e a derivada pode não ser uma densidade no sentido
da identidade \eqref{eq:densidade-integral}. Para os fins deste curso, porém, a
fórmula \eqref{eq:densidade-derivada} é suficiente, e dificilmente encontraremos
problemas ao aplicá-la para determinar a densidade de uma variável aleatória
contínua.

As densidades desempenham para as variáveis contínuas um papel análogo ao
das funções de probabilidade para as variáveis discretas, e não é
surpreendente que as propriedades gerais de umas e de outras sejam muito
parecidas.

\begin{property}
  Se $X$ é uma variável aleatória contínua com densidade $f_X$, então
  \[
    f_X(x) \geqslant 0 \quad \text{para todo } x \in \mathbb{R},
    \qquad
    \int_{-\infty}^{\infty} f_X(x)\,dx = 1.
  \]
\end{property}

Essas duas propriedades são as análogas de $p_X(x) \geqslant 0$ e
$\sum_x p_X(x) = 1$ no caso discreto. A analogia, no entanto, pode ser
perigosa, pois $f_X(x)$ \emph{não} é uma probabilidade e pode muito bem
assumir valores maiores que $1$. No exemplo exponencial, por exemplo, se
$\lambda > 1$, então $f_X$ assume valores maiores que $1$ na vizinhança de
$x = 0$, ainda que $\mathbb{P}(X = 0) = 0$.

O que $f_X(x)$ mede é a densidade de probabilidade nas proximidades de $x$.
Se, além disso, $f_X$ é contínua em $x$ e $\delta x$ é pequeno e positivo,
então, grosso modo, a probabilidade de
$X$ estar ``perto'' de $x$ satisfaz
\begin{align}
  \mathbb{P}(x < X \leqslant x + \delta x)
    &= F_X(x + \delta x) - F_X(x)
       \quad \text{por \eqref{eq:intervalo-fd}} \notag \\[0.2cm]
    &= \int_x^{x + \delta x} f_X(u)\,du
       \quad \text{por \eqref{eq:densidade-integral}} \notag \\[0.2cm]
    &\approx f_X(x)\,\delta x
       \quad \text{para } \delta x \text{ pequeno}.
       \label{eq:densidade-aprox}
\end{align}
A última aproximação usa a continuidade de $f_X$ em $x$: em um intervalo curto,
a área sob $f_X$ é aproximadamente a área de um retângulo de base $\delta x$ e
altura $f_X(x)$.

Portanto, a analogia correta não é entre a densidade $f_X(x)$ e a função de
probabilidade $p_X(x)$, mas sim entre $f_X(x)\,\delta x$ e $p_X(x)$. Essa
leitura fica clara ao compararmos a normalização da densidade com a da
função de probabilidade: os valores de uma função de probabilidade são
substituídos por $f_X(x)\,\delta x$, e a soma, que no caso discreto percorre
apenas os valores positivos (necessariamente em quantidade enumerável), é
substituída pela integral. Uma diferença marcante entre os dois casos é dada
pela primeira parte do teorema a seguir.

\begin{theorem}[Propriedades da densidade]\label{thm:propriedades-densidade}
  Se $X$ é contínua com densidade $f_X$, então
  \begin{equation}\label{eq:prob-ponto}
    \mathbb{P}(X = x) = 0 \quad \text{para todo } x \in \mathbb{R},
  \end{equation}
  e
  \begin{equation}\label{eq:prob-intervalo}
    \mathbb{P}(a \leqslant X \leqslant b) = \int_a^b f_X(u)\,du
    \quad \text{para } a, b \in \mathbb{R} \text{ com } a \leqslant b.
  \end{equation}
\end{theorem}

\begin{proof}
  Vamos provar primeiro \eqref{eq:prob-ponto}. Para isso, observamos que os
  eventos $\{x - \epsilon < X \leqslant x\}$ diminuem para $\{X = x\}$
  quando $\epsilon$ decresce para $0$, de modo que, pela continuidade da
  probabilidade,
  \[
    \mathbb{P}(X = x) = \lim_{\epsilon \downarrow 0}
      \mathbb{P}(x - \epsilon < X \leqslant x).
  \]
  A partir daí, usando sucessivamente \eqref{eq:intervalo-fd} e
  \eqref{eq:densidade-integral}, obtemos
  \[
    \begin{aligned}
      \mathbb{P}(X = x)
        &= \lim_{\epsilon \downarrow 0}
           \bigl[F_X(x) - F_X(x - \epsilon)\bigr] \\
        &= \lim_{\epsilon \downarrow 0} \int_{x - \epsilon}^{x} f_X(u)\,du \\
        &= 0.
    \end{aligned}
  \]
  A última igualdade segue do fato de que a integral sobre um intervalo de
  comprimento $\epsilon$ tende a $0$ quando $\epsilon \downarrow 0$.

  Para a segunda parte, seja $a \leqslant b$. Como $\{X = a\}$ tem
  probabilidade $0$ por \eqref{eq:prob-ponto}, os eventos $\{a \leqslant X
  \leqslant b\}$ e $\{a < X \leqslant b\}$ têm a mesma probabilidade. Logo,
  por \eqref{eq:intervalo-fd} e \eqref{eq:densidade-integral},
  \[
    \begin{aligned}
      \mathbb{P}(a \leqslant X \leqslant b)
        &= \mathbb{P}(a < X \leqslant b) \\
        &= F_X(b) - F_X(a) \\
        &= \int_{-\infty}^{b} f_X(u)\,du - \int_{-\infty}^{a} f_X(u)\,du \\
        &= \int_a^b f_X(u)\,du.
    \end{aligned}
  \]
\end{proof}

O teorema acima mostra que, em uma variável aleatória contínua, a
probabilidade de qualquer valor isolado é nula. Consequentemente, para
variáveis contínuas, incluir ou excluir os extremos de um intervalo não
altera sua probabilidade:
\[
  \mathbb{P}(a \leqslant X \leqslant b) = \mathbb{P}(a < X \leqslant b)
  = \mathbb{P}(a \leqslant X < b) = \mathbb{P}(a < X < b).
\]
Isso contrasta com o caso discreto, em que um único ponto pode concentrar
toda a probabilidade.

\medskip

Recapitulando o que vimos até aqui: toda variável aleatória possui uma
função de distribuição, conforme a \Cref{def:funcao-distribuicao}. No caso
discreto, essa informação é codificada por uma função de probabilidade
$p_X$; no caso contínuo, por uma densidade $f_X$. Essas duas classes não
esgotam, contudo, todas as possibilidades: existem variáveis aleatórias que
não são nem discretas nem contínuas, ou seja, cuja função de distribuição
não é uma escada nem pode ser escrita como a integral de uma densidade.
Tais casos, embora existam, ficam fora do escopo destas notas.

\section{Algumas Funções Densidade Usuais}
\label{sec:densidades-usuais}

Nas seções anteriores vimos que uma variável aleatória contínua $X$ fica
descrita por sua função de distribuição $F_X$ e que, quando $F_X$ é
derivável, sua derivada $f_X$ é a função densidade de probabilidade (fdp)
associada, conforme a \Cref{def:va-continua}. Cabe agora perguntar o
contrário: dada uma função $f$, quando ela é a densidade de alguma variável
aleatória? A resposta é simples e vale a pena enunciá-la antes de listarmos
os exemplos que ocorrem com mais frequência.

Se $f:\mathbb{R}\to[0,\infty)$ é mensurável e satisfaz
\begin{equation}
  f(x) \geqslant 0 \qquad \text{para todo } x \in \mathbb{R}
  \label{eq:densidade-nao-negativa}
\end{equation}
e
\begin{equation}
  \int_{-\infty}^{\infty} f(x)\,dx = 1,
  \label{eq:densidade-normalizacao}
\end{equation}
então $f$ é a função densidade de probabilidade de alguma variável
aleatória. Para confirmar isso, basta definir
\[
  F(x) = \int_{-\infty}^{x} f(u)\,du
\]
e verificar que $F$ satisfaz as três propriedades reunidas na
\Cref{prop:propriedades-funcao-distribuicao}, a saber, $F$ é não decrescente,
contínua à direita, $\lim_{x\to-\infty}F(x)=0$ e
$\lim_{x\to\infty}F(x)=1$. Pela caracterização da
\Cref{rem:caracterizacao-fd}, existe uma variável aleatória $X$ cuja função de
distribuição é $F$. Como $F(x)=\int_{-\infty}^{x}f(u)\,du$, concluímos que $f$
é uma densidade de $X$.

Mostramos abaixo as densidades mais usadas, cada uma acompanhada de seus
parâmetros e de sua região de positividade. Em todos os casos a não negatividade
\eqref{eq:densidade-nao-negativa} é imediata, e a verificação da
normalização \eqref{eq:densidade-normalizacao} é o ponto que merece atenção.
Resumimos a lista em forma de tabela no final desta seção
(\Cref{tab:densidades-usuais}).

\medskip
\noindent\textbf{Distribuição uniforme.}

Dizemos que $X$ tem \emph{distribuição uniforme} no intervalo $(a,b)$, com
$a<b$, quando sua densidade é
\begin{equation}
  f(x) = \begin{cases}
    \dfrac{1}{b-a}, & \text{se } a < x < b, \\[0.3cm]
    0, & \text{caso contrário.}
  \end{cases}
  \label{eq:dens-uniforme}
\end{equation}
Aqui a integral é elementar:
\[
  \int_{-\infty}^{\infty} f(x)\,dx = \int_a^b \frac{1}{b-a}\,dx
  = \frac{b-a}{b-a} = 1.
\]
A densidade é constante sobre o intervalo: não há posição privilegiada
dentro de $(a,b)$, o que justifica o nome.

\medskip
\noindent\textbf{Distribuição exponencial.}

Dizemos que $X$ tem \emph{distribuição exponencial} de parâmetro
$\lambda>0$ quando
\begin{equation}
  f(x) = \begin{cases}
    \lambda e^{-\lambda x}, & \text{se } x > 0, \\[0.2cm]
    0, & \text{se } x \leqslant 0.
  \end{cases}
  \label{eq:dens-exponencial}
\end{equation}
Novamente a normalização sai de uma primitiva imediata, pois
$\frac{d}{dx}\bigl(-e^{-\lambda x}\bigr) = \lambda e^{-\lambda x}$ e
$e^{-\lambda x}\to 0$ quando $x\to\infty$:
\[
  \int_{-\infty}^{\infty} f(x)\,dx
  = \int_0^{\infty} \lambda e^{-\lambda x}\,dx
  = \Bigl[-e^{-\lambda x}\Bigr]_{0}^{\infty} = 1.
\]

\medskip
\noindent\textbf{Distribuição normal (gaussiana).}

Dizemos que $X$ tem \emph{distribuição normal} (ou \emph{gaussiana}) de
parâmetros $\mu\in\mathbb{R}$ e $\sigma>0$, e escrevemos
$X\sim \mathrm{N}(\mu,\sigma^2)$, quando
\begin{equation}
  f(x) = \frac{1}{\sqrt{2\pi\sigma^2}}
  \exp\!\left(-\frac{1}{2\sigma^2}(x-\mu)^2\right),
  \qquad x \in \mathbb{R}.
  \label{eq:dens-normal}
\end{equation}
Para verificar que a integral vale $1$, usamos a integral gaussiana
\begin{equation}
  \int_{-\infty}^{\infty} e^{-x^2}\,dx = \sqrt{\pi}.
  \label{eq:gaussiana-integral}
\end{equation}
Com efeito, a mudança de variáveis $x = \mu + \sigma\sqrt{2}\,z$ dá
$dx = \sigma\sqrt{2}\,dz$ e
\[
  \int_{-\infty}^{\infty} e^{-\frac{1}{2\sigma^2}(x-\mu)^2}\,dx
  = \sigma\sqrt{2}\int_{-\infty}^{\infty} e^{-z^2}\,dz
  = \sigma\sqrt{2\pi},
\]
onde usamos \eqref{eq:gaussiana-integral}. Dividindo pelo fator
$\sqrt{2\pi\sigma^2}$ obtemos exatamente $1$. A constante de normalização
foi escolhida justamente para isso.

\medskip
\noindent\textbf{Distribuição de Cauchy.}

Dizemos que $X$ tem \emph{distribuição de Cauchy} quando
\begin{equation}
  f(x) = \frac{1}{\pi(1+x^2)}, \qquad x \in \mathbb{R}.
  \label{eq:dens-cauchy}
\end{equation}
A normalização segue da derivada da função arco-tangente, pois
$\frac{d}{dx}\arctan x = \frac{1}{1+x^2}$ e
$\lim_{x\to\pm\infty}\arctan x = \pm\frac{\pi}{2}$:
\[
  \int_{-\infty}^{\infty} \frac{1}{\pi(1+x^2)}\,dx
  = \frac{1}{\pi}\Bigl[\arctan x\Bigr]_{-\infty}^{\infty}
  = \frac{1}{\pi}\left(\frac{\pi}{2}+\frac{\pi}{2}\right) = 1.
\]

\medskip
\noindent\textbf{Distribuição gama.}

Dizemos que $X$ tem \emph{distribuição gama} de parâmetros $w>0$ e
$\lambda>0$ quando
\begin{equation}
  f(x) = \begin{cases}
    \dfrac{1}{\Gamma(w)}\,\lambda^w x^{w-1} e^{-\lambda x}, & \text{se } x > 0, \\[0.3cm]
    0, & \text{se } x \leqslant 0,
  \end{cases}
  \label{eq:dens-gama}
\end{equation}
onde $\Gamma$ é a \emph{função gama}, definida por
\begin{equation}
  \Gamma(w) = \int_0^{\infty} x^{w-1} e^{-x}\,dx.
  \label{eq:dens-gama-func}
\end{equation}
A integral em \eqref{eq:dens-gama-func} converge para todo $w>0$: perto de
$0$, o integrando é comparável a $x^{w-1}$, cuja integral converge para
$w>0$, enquanto, no infinito, o fator $e^{-x}$ garante a convergência. Para
checar a normalização de \eqref{eq:dens-gama}, fazemos a mudança
$y=\lambda x$, de modo que $dx = dy/\lambda$:
\[
  \int_{-\infty}^{\infty} f(x)\,dx
  = \frac{\lambda^w}{\Gamma(w)}\int_0^{\infty} x^{w-1} e^{-\lambda x}\,dx
  = \frac{1}{\Gamma(w)}\int_0^{\infty} y^{w-1} e^{-y}\,dy
  = \frac{\Gamma(w)}{\Gamma(w)} = 1.
\]
Ou seja, a distribuição gama integra $1$ \emph{por definição} de $\Gamma$.

Observamos ainda que, para inteiro positivo $w=n$, a função gama
recupera o fatorial. De fato, integrando por partes,
$\Gamma(w) = (w-1)\Gamma(w-1)$ para $w>1$, e como $\Gamma(1)=1$ segue que
\[
  \Gamma(n) = (n-1)! \qquad \text{para } n=1,2,3,\dots
\]
Este é o conteúdo do \Cref{exercise:funcao-gama-recorrencia}, ao qual
remetemos o leitor para os detalhes da indução.

\medskip
\noindent\textbf{Distribuição beta.}

Dizemos que $X$ tem \emph{distribuição beta} de parâmetros $s>0$ e $t>0$
quando
\begin{equation}
  f(x) =
  \begin{cases}
    \dfrac{1}{B(s,t)}\,x^{s-1}(1-x)^{t-1}, & \text{se } 0 < x < 1, \\[0.3cm]
    0, & \text{caso contrário},
  \end{cases}
  \label{eq:dens-beta}
\end{equation}
onde $B$ é a \emph{função beta},
\begin{equation}
  B(s,t) = \int_0^1 x^{s-1}(1-x)^{t-1}\,dx.
  \label{eq:dens-beta-func}
\end{equation}
Aqui a normalização é tautológica: a constante $B(s,t)$ foi escolhida
exatamente para que a integral de $f$ em \eqref{eq:dens-beta} seja
\[
  \frac{1}{B(s,t)}\int_0^1 x^{s-1}(1-x)^{t-1}\,dx
  = \frac{B(s,t)}{B(s,t)} = 1.
\]
No caso particular $s=t=1$, a densidade vale $1$ em $(0,1)$ e $0$ fora
desse intervalo. Os valores nos extremos são irrelevantes, e obtemos a
distribuição uniforme em $[0,1]$, o que concorda com
\eqref{eq:dens-uniforme} tomando $a=0$ e $b=1$.

\medskip
\noindent\textbf{Distribuição qui-quadrado.}

Dizemos que $X$ tem \emph{distribuição qui-quadrado com $n$ graus de
liberdade}, e escrevemos $X\sim\chi_n^2$, quando
\begin{equation}
  f(x) = \begin{cases}
    \dfrac{1}{2\Gamma\!\bigl(\frac{n}{2}\bigr)}
    \Bigl(\dfrac{x}{2}\Bigr)^{\frac{n}{2}-1} e^{-\frac{x}{2}}, & \text{se } x > 0, \\[0.4cm]
    0, & \text{se } x \leqslant 0.
  \end{cases}
  \label{eq:dens-qui-quadrado}
\end{equation}
Comparando \eqref{eq:dens-qui-quadrado} com \eqref{eq:dens-gama}, vemos que
esta é precisamente a distribuição gama com parâmetros $w=\frac{n}{2}$ e
$\lambda=\frac{1}{2}$. Logo a normalização já foi verificada no caso gama, e
listamos a distribuição separadamente apenas porque ela aparece com muita
frequência em estatística.

\medskip

A tabela abaixo reúne as densidades discutidas acima. A coluna final registra a
região de positividade, isto é, o conjunto onde $f$ é estritamente positiva.
Ela não deve ser confundida com o suporte topológico, que também contém pontos
de fronteira.
\begin{table}[H]
  \centering
  \small
  \begin{tabular}{@{}llll@{}}
    \toprule
    Distribuição & Parâmetros & Densidade $f(x)$ & Região de positividade \\
    \midrule
    Uniforme & $a<b$ &
    $\dfrac{1}{b-a}$ & $a<x<b$ \\[0.35cm]
    Exponencial & $\lambda>0$ &
    $\lambda e^{-\lambda x}$ & $x>0$ \\[0.25cm]
    Normal & $\mu\in\mathbb{R}$, $\sigma>0$ &
    $\dfrac{1}{\sqrt{2\pi\sigma^2}}
     e^{-\frac{(x-\mu)^2}{2\sigma^2}}$ & $x\in\mathbb{R}$ \\[0.45cm]
    Cauchy & nenhum &
    $\dfrac{1}{\pi(1+x^2)}$ & $x\in\mathbb{R}$ \\[0.4cm]
    Gama & $w>0$, $\lambda>0$ &
    $\dfrac{1}{\Gamma(w)}\lambda^w x^{w-1}e^{-\lambda x}$ & $x>0$ \\[0.4cm]
    Beta & $s>0$, $t>0$ &
    $\dfrac{1}{B(s,t)}x^{s-1}(1-x)^{t-1}$ & $0<x<1$ \\[0.4cm]
    Qui-quadrado & $n\in\mathbb{N}$ &
    $\dfrac{1}{2\Gamma(\frac{n}{2})}
     \bigl(\frac{x}{2}\bigr)^{\frac{n}{2}-1}e^{-\frac{x}{2}}$ & $x>0$ \\
    \bottomrule
  \end{tabular}
  \caption{As densidades usuais, com seus parâmetros e regiões de positividade. Cada
  constante de normalização foi escolhida para que a integral da densidade
  sobre $\mathbb{R}$ seja igual a $1$.}
  \label{tab:densidades-usuais}
\end{table}

\medskip


A lista apresentada na \Cref{sec:densidades-usuais} é um compêndio de algumas
das densidades mais comuns, e não esperamos que ela, sozinha, seja
inspiradora. É difícil motivar adequadamente cada uma dessas densidades neste
momento, mas vamos precisar retomá-las mais adiante, quando as encontrarmos em
ação.

\medskip

Não é sempre trivial verificar que as funções da lista são realmente densidades. A condição de não negatividade não oferece dificuldade; o cuidado está em conferir que a integral vale $1$. Para a densidade normal, por exemplo, precisamos da integral definida padrão
\[
  \int_{-\infty}^{\infty} e^{-x^2}\,dx = \sqrt{\pi}.
\]
Um esboço da prova dessa identidade, por meio de coordenadas polares, está no \Cref{exercise:integral-gaussiana}.

Os fatores constantes que aparecem nas densidades usuais foram escolhidos apenas para que as funções resultantes tenham integral $1$. Considere, por exemplo, a função
\[
  g(x) = \frac{1}{1 + x^2} \quad \text{para } x \in \mathbb{R}.
\]
Como
\[
  \int_{-\infty}^{\infty} g(x)\,dx = \pi,
\]
$g$ não é uma densidade. No entanto, a função normalizada
\[
  f(x) = \frac{1}{\pi} g(x)
\]
é uma densidade, a qual corresponde à distribuição de Cauchy apresentada na
\Cref{sec:densidades-usuais}.


\medskip

Resta uma relação entre as funções gama e beta que simplifica o cálculo com
a densidade beta e que usaremos mais adiante. Ela expressa $B(s,t)$ em
termos de $\Gamma$.

\begin{proposition}
  \label{prop:beta-gama}
  Para todos $s>0$ e $t>0$ vale
  \[
    B(s,t) = \frac{\Gamma(s)\Gamma(t)}{\Gamma(s+t)}.
  \]
\end{proposition}

\begin{proof}
  Partimos do produto das duas integrais que definem $\Gamma(s)$ e
  $\Gamma(t)$. Como os integrandos são não negativos, podemos escrever o
  produto como uma integral dupla sobre o quadrante $(0,\infty)^2$, sem
  qualquer preocupação com convergência:
  \[
    \Gamma(s)\Gamma(t)
    = \int_0^{\infty}\!\!\int_0^{\infty}
      e^{-(u+v)} u^{s-1} v^{t-1}\,du\,dv.
  \]
  Fazemos agora a mudança de variáveis
  \[
    u = r\cos^2\theta, \qquad v = r\sin^2\theta,
  \]
  com $r>0$ e $\theta\in(0,\pi/2)$. Essa mudança é sugerida pela
  identidade $\cos^2\theta+\sin^2\theta=1$: ela transforma a
  soma $u+v$ na variável radial $r$, de modo que o fator $e^{-(u+v)}$
  passa a depender apenas de $r$. O determinante jacobiano é
  \[
    J = \det
    \begin{pmatrix}
      \cos^2\theta & -2r\sin\theta\cos\theta \\[0.2cm]
      \sin^2\theta & 2r\sin\theta\cos\theta
    \end{pmatrix}
    = 2r\sin\theta\cos\theta,
  \]
  que é positivo no domínio considerado, logo $|J| = J$. Além disso,
  $u+v = r$ e
  \[
    u^{s-1}v^{t-1}
    = r^{s+t-2}\cos^{2(s-1)}\theta\,\sin^{2(t-1)}\theta.
  \]
  Substituindo na integral dupla, obtemos
  \begin{align*}
    \Gamma(s)\Gamma(t)
    &= \int_0^{\pi/2}\!\!\int_0^{\infty}
       e^{-r} r^{s+t-2}\cos^{2(s-1)}\theta\,\sin^{2(t-1)}\theta\,
       2r\sin\theta\cos\theta\,dr\,d\theta \\
    &= \left(\int_0^{\infty} e^{-r} r^{s+t-1}\,dr\right)
       \left(2\int_0^{\pi/2}
       \cos^{2s-1}\theta\,\sin^{2t-1}\theta\,d\theta\right),
  \end{align*}
  onde separamos as variáveis porque o domínio é um produto e o integrando
  é um produto de um fator que depende só de $r$ por outro que depende só
  de $\theta$. A integral em $r$ é exatamente a definição de
  $\Gamma(s+t)$.

  Resta identificar a integral em $\theta$ com $B(s,t)$. Para isso,
  pomos $x=\cos^2\theta$. Então $dx = -2\cos\theta\sin\theta\,d\theta$
  e, quando $\theta$ vai de $0$ a $\pi/2$, $x$ vai de $1$ a $0$. Escrevendo
  \[
    \cos^{2s-1}\theta = x^{s-1}\cos\theta,
    \qquad
    \sin^{2t-1}\theta = (1-x)^{t-1}\sin\theta,
  \]
  segue que
  \[
    \begin{aligned}
      2\int_0^{\pi/2}
      \cos^{2s-1}\theta\,\sin^{2t-1}\theta\,d\theta
      &= 2\int_1^0 x^{s-1}(1-x)^{t-1}
        \cos\theta\sin\theta\,
        \frac{-dx}{2\cos\theta\sin\theta} \\[0.2cm]
      &= \int_0^1 x^{s-1}(1-x)^{t-1}\,dx,
    \end{aligned}
  \]
  que é $B(s,t)$ por \eqref{eq:dens-beta-func}. Juntando os dois fatores,
  \[
    \Gamma(s)\Gamma(t) = \Gamma(s+t)\,B(s,t),
  \]
  e como $\Gamma(s+t)>0$, dividindo por esse valor chegamos a
  $B(s,t) = \dfrac{\Gamma(s)\Gamma(t)}{\Gamma(s+t)}$, como queríamos.
\end{proof}

A \Cref{prop:beta-gama} mostra, em particular, que a constante de
normalização da densidade beta também pode ser escrita em termos da função
gama: $B(s,t) = \Gamma(s)\Gamma(t)/\Gamma(s+t)$. Com isso encerramos a
lista das densidades usuais, que retomaremos quando essas distribuições
aparecerem em situações concretas.

\section{Exercícios propostos}
\label{sec:exercicios-funcoes-distribuicao}

Os exercícios a seguir retomam os conceitos e exemplos desenvolvidos ao longo
da nota: funções de distribuição, densidades e as distribuições usuais.
Salvo indicação em contrário, $X$ denota uma variável aleatória em um espaço
de probabilidade fixado.

\medskip
\noindent\textbf{Funções de distribuição}

\begin{exercise}\label{exercise:soma-massas-intervalo}
  Seja $X$ uma variável aleatória que assume valores inteiros e satisfaz
  $\mathbb{P}(X = k) = p_k$ para $k = \dots, -1, 0, 1, \dots$. Mostre que a
  função de distribuição de $X$ satisfaz
  \[
    F_X(b) - F_X(a) = p_{a+1} + p_{a+2} + \dots + p_b
  \]
  para todos os inteiros $a, b$ com $a < b$.
\end{exercise}

\begin{exercise}\label{exercise:descontinuidade-ponto-massa}
  Seja $X$ uma variável aleatória e seja $c$ um número real tal que
  $\mathbb{P}(X = c) > 0$. Mostre que a função de distribuição $F_X(x)$ de $X$
  é descontínua no ponto $x = c$. A recíproca é verdadeira?
\end{exercise}

\begin{exercise}\label{exercise:fd-maximo-zero}
  Expresse a função de distribuição de $Y = \max\{0, X\}$ em termos da função
  de distribuição $F_X$ de $X$.
\end{exercise}

\begin{exercise}\label{exercise:existencia-mediana}
  Dizemos que o número real $m$ é uma \emph{mediana} da variável aleatória $X$
  se
  \[
    \mathbb{P}(X < m) \leqslant \frac{1}{2} \leqslant \mathbb{P}(X \leqslant m).
  \]
  Mostre que toda variável aleatória possui pelo menos uma mediana.
\end{exercise}

\medskip
\noindent\textbf{Exemplos de funções de distribuição}

\begin{exercise}\label{exercise:mistura-fd}
  Mostre que se $F_1$ e $F_2$ são funções de distribuição, então também é função
  de distribuição a função
  \[
    F(x) = \alpha F_1(x) + (1 - \alpha) F_2(x)
  \]
  para qualquer $\alpha$ satisfazendo $0 \leqslant \alpha \leqslant 1$.
\end{exercise}

\begin{exercise}\label{exercise:normalizacao-exponencial-modulo}
  Seja
  \[
    F(x) = c \int_{-\infty}^{x} e^{-|u|}\,du \qquad \text{para } x \in
    \mathbb{R}.
  \]
  Para qual valor de $c$ a função $F$ é uma função de distribuição?
\end{exercise}

\medskip
\noindent\textbf{Variáveis aleatórias contínuas}

\begin{exercise}
\label{exercise:densidade-linear-crescente}
  Uma variável aleatória \(X\) tem função densidade de probabilidade
  \[
    f(x) = \begin{cases} 2x & \text{se } 0 < x < 1, \\ 0 & \text{caso
      contrário.} \end{cases}
  \]
  Determine a função de distribuição de \(X\).
\end{exercise}

\begin{exercise}
\label{exercise:densidade-exponencial-bilateral}
  Se \(X\) tem função densidade de probabilidade
  \[
    f(x) = \frac{1}{2}e^{-|x|} \qquad \text{para } x \in \mathbb{R},
  \]
  determine a função de distribuição de \(X\). Esta distribuição é chamada
  de distribuição \emph{exponencial bilateral} (ou \emph{dupla}).
\end{exercise}

\begin{exercise}
\label{exercise:fd-por-partes}
  Se \(X\) tem função de distribuição
  \[
    F(x) = \begin{cases} \frac{1}{2(1+x^2)} & \text{para } -\infty < x
      \leqslant 0, \\[0.2cm] \frac{1+2x^2}{2(1+x^2)} & \text{para } 0 < x <
      \infty, \end{cases}
  \]
  mostre que \(X\) é contínua e determine a sua função densidade de
  probabilidade.
\end{exercise}

\begin{exercise}
\label{exercise:densidade-valor-extremo}
  Determine a função de distribuição da chamada densidade de ``valor
  extremo''
  \[
    f(x) = \exp(-x - e^{-x}) \qquad \text{para } x \in \mathbb{R}.
  \]
\end{exercise}

\medskip
\noindent\textbf{Densidades usuais}

\begin{exercise}
  \label{exercise:gama-exponencial}
  Para quais valores dos seus parâmetros a distribuição gama é também uma distribuição exponencial?
\end{exercise}

\begin{exercise}
  \label{exercise:funcao-gama-recorrencia}
  Mostre que a função gama $\Gamma(w)$ satisfaz
  \[
    \Gamma(w) = (w - 1)\Gamma(w - 1)
  \]
  para $w > 1$, e deduza que
  \[
    \Gamma(n) = (n - 1)!
  \]
  para $n = 1, 2, 3, \ldots$.
\end{exercise}

\begin{exercise}
  \label{exercise:integral-gaussiana}
  Seja
  \[
    I = \int_{-\infty}^{\infty} e^{-x^2}\,dx.
  \]
  Passando a coordenadas polares, mostre que
  \[
    I^2 = \iint_{\mathbb{R}^2} e^{-x^2 - y^2}\,dx\,dy
        = \int_{\theta=0}^{2\pi} \int_{r=0}^{\infty} e^{-r^2} r\,dr\,d\theta,
  \]
  e deduza que $I = \sqrt{\pi}$.
\end{exercise}

\begin{exercise}
  \label{exercise:densidade-arcoseno}
  Mostre que a função densidade
  \[
    f(x) = \begin{cases}
      \dfrac{1}{\pi\sqrt{x(1 - x)}} & \text{se } 0 < x < 1, \\[0.3cm]
      0 & \text{caso contrário},
    \end{cases}
  \]
  tem função de distribuição da forma
  \[
    F(x) = c\operatorname{arcsen}\sqrt{x} \qquad \text{se } 0 < x < 1,
  \]
  e determine a constante $c$.
\end{exercise}

\bibliographystyle{plain}
\bibliography{../arquivos-aux/referencias}

\end{document}
